% Generated draft: requires mathematical and source review.
\documentclass[11pt]{article}
\usepackage[margin=25mm]{geometry}
\usepackage{fontspec}
\setmainfont{Latin Modern Roman}
\usepackage{amsmath,amssymb,mathtools,unicode-math}
\setmathfont{Latin Modern Math}
\usepackage{enumitem,needspace,xurl,hyperref}
\setlength{\parindent}{0pt}
\setlength{\parskip}{5pt}
\setlength{\emergencystretch}{3em}
\newcommand{\N}{\mathbb{N}}
\newcommand{\Z}{\mathbb{Z}}
\newcommand{\Q}{\mathbb{Q}}
\newcommand{\R}{\mathbb{R}}
\newcommand{\C}{\mathbb{C}}
\newcommand{\F}{\mathbb{F}}
\newcommand{\A}{\mathbb{A}}
\newcommand{\PP}{\mathbb{P}}
\newcommand{\E}{\mathbb{E}}
\newcommand{\eps}{\varepsilon}
\newcommand{\dd}{\,\mathrm{d}}
\begin{document}
\section*{UnsolvedMath ID: 120}
\textbf{Source problem code:} catalogue entry
\subsection{Definitions and mathematical statement}
For an integer $m\geq 2$, let $P^{+}(m)$ denote its greatest prime factor, and put $P^{+}(1)=1$. A positive integer $m$ is $y$-smooth when $P^{+}(m)\leq y$. The supplied source asks whether every $n\leq N$ is a sum of two integers whose prime factors are at most $N^{\varepsilon}$. Its asymptotic interpretation requires a range for $n$ and a convention concerning the summands; these are not explicitly supplied.

Under the usual interpretation with positive summands, define, for $n\geq2$,
\[
f(n)=\min_{\substack{a,b\in\Z_{\geq1}\\a+b=n}}\max\{P^{+}(a),P^{+}(b)\}.
\]
The target identified in the supplied dated literature assessment is
\[
\forall\varepsilon>0\ \exists n_0(\varepsilon)\ \forall n\in\Z_{\geq n_0(\varepsilon)}:\quad f(n)\leq n^{\varepsilon},
\]
that is, $f(n)=n^{o(1)}$. Equivalently, with the same positive-summand convention, the uniform formulation is
\[
\forall\varepsilon>0\ \exists N_0(\varepsilon)\ \forall N\in\Z_{\geq N_0(\varepsilon)}\ \forall n\in\Z,\ 2\leq n\leq N:\quad
\exists a,b\in\Z_{\geq1},\quad a+b=n,\quad P^{+}(a),P^{+}(b)\leq N^{\varepsilon}.
\]
The lower cutoff $n\geq2$, positivity, and sufficiently large $N$ make explicit conventions needed for this formulation; the unqualified wording ``every $n\leq N$'' does not specify them. In particular $n=1$ cannot be represented by two positive integers. If zero or signed summands were intended, their permitted domains and smoothness convention would need to be specified separately.

The supplied literature triage, checked 2026-08-17, labels the problem partially solved and reports Balog's estimate
\[
f(n)\ll_{\varepsilon}n^{4/(9\sqrt e)+\varepsilon}
\quad(\varepsilon>0),
\]
with exponent $4/(9\sqrt e)\approx0.2695$. It states that arbitrarily small positive exponents remain open and describes a least-quadratic-nonresidue consequence and a Burgess-type barrier, without supplying their precise formulations. These are dated source assessments, not a new verification. The registry's primary status remains ``open''. The source identifies the local code catalogue entry with Green Problem 59 and Erd\H{o}s Problem 334 by statement.
\subsection{Short English statement}
Does every sufficiently large integer split into two positive integers whose largest prime factors are smaller than every prescribed positive power of that integer? The source also gives a uniform formulation for $n\leq N$, but does not explicitly state the range endpoints or summand convention. Its dated assessment reports partial progress with exponent about $0.2695$.
\subsection{Sources}
\begin{itemize}
\item UnsolvedMath local registry, permanent ID 120, source code catalogue entry, ``Sums of Smooth Numbers'', Ben Green's 100 Open Problems collection. Dataset attribution: Ulam AI and the UnsolvedMath Contributors; curation and metadata under CC BY 4.0.
\item T. F. Bloom, Erd\H{o}s Problem 334, \url{https://www.erdosproblems.com/334}. The supplied triage records a check on 2026-08-17 and cites its open label and Balog bound.
\item B. Green, \emph{100 Open Problems}, Problem 59, version recorded as retrieved 2026-08-17, \url{https://people.maths.ox.ac.uk/greenbj/papers/open-problems.pdf}. The supplied triage cites the same exponent and discussion of least quadratic nonresidues.
\item The supplied record attributes the partial bound to Balog but gives no direct bibliographic citation to Balog's work.
\end{itemize}
\end{document}
